% textmakro.sty - Dokumentation
% Autor: OTS
% Datum: 28. Februar 2026

\documentclass[12pt,a4paper]{scrartcl}

% Pakete
\usepackage[ngerman]{babel}
\usepackage[T1]{fontenc}
\usepackage{lmodern}
\usepackage{geometry}
\usepackage{xcolor}
\usepackage{tcolorbox}
\usepackage{booktabs}
\usepackage{longtable}
\usepackage{listings}
\usepackage{textmakro}

% Seitengeometrie
\geometry{
  left=2.5cm,
  right=2.5cm,
  top=2.5cm,
  bottom=2.5cm
}

% Listings-Konfiguration
\lstset{
  basicstyle=\ttfamily\small,
  breaklines=true,
  frame=single,
  backgroundcolor=\color{gray!10},
  keywordstyle=\color{blue},
  commentstyle=\color{green!60!black},
  stringstyle=\color{red},
  showstringspaces=false,
  numbers=left,
  numberstyle=\tiny\color{gray},
  language=[LaTeX]TeX
}

% Titel
\title{\textbf{textmakro.sty}\\[0.5em]\Large Schnellreferenz}
\author{OTS}
\date{28. Februar 2026}

\begin{document}

\maketitle
\tableofcontents
\newpage

%%%%%%%%%%%%%%%%%%%%%%%%%%%%%%%%%%%%%%%%%%%%%%%%%%%%%%%%%
\section{Installation}
%%%%%%%%%%%%%%%%%%%%%%%%%%%%%%%%%%%%%%%%%%%%%%%%%%%%%%%%%

Um das \texttt{textmakro}-Paket zu verwenden:

\begin{lstlisting}
\usepackage{textmakro}
\end{lstlisting}

Erfordert: \texttt{colors.sty} und \texttt{xparse}

%%%%%%%%%%%%%%%%%%%%%%%%%%%%%%%%%%%%%%%%%%%%%%%%%%%%%%%%%
\section{Verwendung}
%%%%%%%%%%%%%%%%%%%%%%%%%%%%%%%%%%%%%%%%%%%%%%%%%%%%%%%%%

Zwei Modi:
\begin{itemize}
  \item \textbf{Textfarbe:} \lstinline|\makro{Text}|
  \item \textbf{Box:} \lstinline|\makro[b]{Text}|
\end{itemize}

%%%%%%%%%%%%%%%%%%%%%%%%%%%%%%%%%%%%%%%%%%%%%%%%%%%%%%%%%
\section{Makro-{\"U}bersicht}
%%%%%%%%%%%%%%%%%%%%%%%%%%%%%%%%%%%%%%%%%%%%%%%%%%%%%%%%%

%--------------------------------------------------------
\subsection{Verben}
%--------------------------------------------------------

\begin{tabular}{lll}
\toprule
\textbf{Makro} & \textbf{Farbe} & \textbf{Verwendung} \\
\midrule
\lstinline|\verb{Text}| & Gr{\"u}n & Normale Verben \\
\lstinline|\verbI{Text}| & Gelb & Imperative \\
\lstinline|\verbP{Text}| & Gr{\"u}n + P & Passive \\
\bottomrule
\end{tabular}

Beispiele:
\begin{lstlisting}
Abraham \verb{ging} nach Ur.
\verbI{Gehe} und \verbI[b]{verkuende}!
\end{lstlisting}

%--------------------------------------------------------
\subsection{Nomen}
%--------------------------------------------------------

\begin{tabular}{lll}
\toprule
\textbf{Makro} & \textbf{Farbe} & \textbf{Verwendung} \\
\midrule
\lstinline|\nomen{Text}| & Orange & Substantive \\
\bottomrule
\end{tabular}

Beispiel:
\begin{lstlisting}
Das \nomen{Haus} des Herrn.
\end{lstlisting}

%--------------------------------------------------------
\subsection{Konjunktionen}
%--------------------------------------------------------

\begin{tabular}{lll}
\toprule
\textbf{Makro} & \textbf{Farbe} & \textbf{Verwendung} \\
\midrule
\lstinline|\konj{Text}| & Rot & Bindew{\"o}rter \\
\bottomrule
\end{tabular}

Beispiel:
\begin{lstlisting}
Abraham \konj{und} Sara.
\end{lstlisting}

%--------------------------------------------------------
\subsection{Orte}
%--------------------------------------------------------

\begin{tabular}{lll}
\toprule
\textbf{Makro} & \textbf{Farbe} & \textbf{Verwendung} \\
\midrule
\lstinline|\ort{Text}| & Magenta & Geografische Orte \\
\bottomrule
\end{tabular}

Beispiel:
\begin{lstlisting}
Von \ort{Jerusalem} nach \ort[b]{Bethlehem}.
\end{lstlisting}

%--------------------------------------------------------
\subsection{Personen}
%--------------------------------------------------------

\begin{tabular}{lll}
\toprule
\textbf{Makro} & \textbf{Farbe} & \textbf{Verwendung} \\
\midrule
\lstinline|\person{Text}| & Blau & Personen \\
\lstinline|\mensch{Text}| & Blau & Menschen \\
\bottomrule
\end{tabular}

Beispiel:
\begin{lstlisting}
\person{Abraham} war ein \mensch{Mensch} des Glaubens.
\end{lstlisting}

%--------------------------------------------------------
\subsection{Theologische Personen}
%--------------------------------------------------------

\begin{tabular}{lll}
\toprule
\textbf{Makro} & \textbf{Farbe} & \textbf{Verwendung} \\
\midrule
\lstinline|\gott{Text}| & Himmelsblau & Gott \\
\lstinline|\jesus{Text}| & Himmelsblau & Jesus \\
\lstinline|\geist{Text}| & Himmelsblau & Heiliger Geist \\
\lstinline|\israel{Text}| & Lila & Israel \\
\lstinline|\teufel{Text}| & Rot (hell) & Teufel \\
\bottomrule
\end{tabular}

Beispiel:
\begin{lstlisting}
\gott{Gott} sandte \jesus[b]{Jesus}.
\end{lstlisting}

%--------------------------------------------------------
\subsection{Hilfsmakros}
%--------------------------------------------------------

\begin{tabular}{ll}
\toprule
\textbf{Makro} & \textbf{Beschreibung} \\
\midrule
\lstinline|\annot{Text}| & Annotation (grau, klein, kursiv) \\
\bottomrule
\end{tabular}

%%%%%%%%%%%%%%%%%%%%%%%%%%%%%%%%%%%%%%%%%%%%%%%%%%%%%%%%%
\section{Farb{\"u}bersicht}
%%%%%%%%%%%%%%%%%%%%%%%%%%%%%%%%%%%%%%%%%%%%%%%%%%%%%%%%%

\begin{tabular}{lll}
\toprule
\textbf{Kategorie} & \textbf{Farbe} & \textbf{RGB} \\
\midrule
Verben & Gr{\"u}n & 50,205,50 \\
Imperativ & Gelb & 255,255,0 \\
Nomen & Orange & 255,140,0 \\
Konjunktionen & Rot & 255,0,0 \\
Orte & Magenta & 255,0,255 \\
Personen & Blau & 100,100,205 \\
Gott/Jesus/Geist & Himmelsblau & 180,210,255 \\
Israel & Lila & 220,180,255 \\
Teufel & Rot (hell) & 255,180,180 \\
\bottomrule
\end{tabular}

%%%%%%%%%%%%%%%%%%%%%%%%%%%%%%%%%%%%%%%%%%%%%%%%%%%%%%%%%
\section{Boolean-Schalter}
%%%%%%%%%%%%%%%%%%%%%%%%%%%%%%%%%%%%%%%%%%%%%%%%%%%%%%%%%

Farben k{\"o}nnen deaktiviert werden:

\begin{lstlisting}
\ExplSyntaxOn
\bool_set_false:N \l_color_verb_bool
\bool_set_false:N \l_color_nomen_bool
...
\ExplSyntaxOff
\end{lstlisting}

\subsection{Verf{\"u}gbare Boolean-Schalter}

\begin{tabular}{ll}
\toprule
\textbf{Variable} & \textbf{Makro} \\
\midrule
\lstinline|\l_color_verb_bool| & \lstinline|\verb{}| \\
\lstinline|\l_color_verbI_bool| & \lstinline|\verbI{}| \\
\lstinline|\l_color_verbP_bool| & \lstinline|\verbP{}| \\
\lstinline|\l_color_nomen_bool| & \lstinline|\nomen{}| \\
\lstinline|\l_color_konj_bool| & \lstinline|\konj{}| \\
\lstinline|\l_color_ort_bool| & \lstinline|\ort{}| \\
\lstinline|\l_color_person_bool| & \lstinline|\person{}, \mensch{}| \\
\lstinline|\l_color_gott_bool| & \lstinline|\gott{}| \\
\lstinline|\l_color_jesus_bool| & \lstinline|\jesus{}| \\
\lstinline|\l_color_geist_bool| & \lstinline|\geist{}| \\
\lstinline|\l_color_israel_bool| & \lstinline|\israel{}| \\
\lstinline|\l_color_teufel_bool| & \lstinline|\teufel{}| \\
\bottomrule
\end{tabular}

%%%%%%%%%%%%%%%%%%%%%%%%%%%%%%%%%%%%%%%%%%%%%%%%%%%%%%%%%
\section{Tipps}
%%%%%%%%%%%%%%%%%%%%%%%%%%%%%%%%%%%%%%%%%%%%%%%%%%%%%%%%%

\begin{enumerate}
  \item Verwenden Sie entweder Textfarben ODER Box-Farben
  \item Zu viele Boxen wirken {\"u}berladen
  \item Boxen f{\"u}r Schl{\"u}sselw{\"o}rter
  \item Boolean-Schalter f{\"u}r Druck nutzen
\end{enumerate}

\end{document}
